\documentclass[11pt,a4paper]{article}
\usepackage[margin=25mm]{geometry}
\usepackage{amsmath,amssymb,amsthm,hyperref}
\newtheorem{lemma}{Lemma}
\newtheorem{proposition}[lemma]{Proposition}
\newcommand{\F}{\mathbb F}
\newcommand{\Om}{\Omega}
\setlength{\emergencystretch}{2em}
\title{Choosing an orthogonal torus element to separate the triality modules}
\author{Human Compression project}\date{}
\begin{document}\maketitle
This draft supplies the triality interface in plus dimension eight.
It awaits distinct review and does not classify all maximal overgroups.

Let $G=P\Om_8^+(q)$, with inner-diagonal automorphism group $J$.
We retain the standard automorphism structure: modulo $J$, the diagram
group $S_3$ commutes with the cyclic field group. The three diagram points
are the natural and two half-spin eight-dimensional projective
representations. The stabilizer of one point acts by natural projective
conformal maps in that representation. The Clifford and triality
descriptions work in characteristic two as well; see~\cite[Sections~3.1
and~3.4]{BT}, in particular Theorem~3.12 and the isomorphisms
(3.20)--(3.21). These are borrowed structural statements.

Put $A=q^3+1$, $B=q+1$, $b=q^2-q+1$, so $A=Bb$.
Choose $g$ of order $A$ in the norm-one torus of
$\F_{q^6}/\F_{q^3}$. On negative spaces of dimensions six and two use
primitive multipliers $g$ and $h=g^{ub}$, with $u$ a unit modulo $B$.
Let $X$ be their direct sum and $x$ its image in $G$.
The accepted torus proof gives $X\in\Om_8^+(q)$,
$|x|=A/\gcd(2,q-1)$, and the connected adjoint centralizer
\[
                      C_J(x)=T,\qquad |T|=AB.                 \tag{1}
\]
In particular its order is not replaced by $\gcd(2,q-1)|x|$.

\begin{lemma}\label{weights}
Modulo scalar signs, the natural eigenvalue exponents relative to $g$
are
\[
                       \pm1,\quad\pm q,\quad\pm q^2,\quad\pm ub.
\]
The two half-spin lists are
\[
 c_\epsilon q^j\ (0\le j<6),\quad d_\epsilon,-d_\epsilon,\qquad
 c_\epsilon=\frac{1+q+q^2+\epsilon ub}{2},\quad
 d_\epsilon=\frac{(1-\epsilon u)b}{2},\quad\epsilon=\pm1.       \tag{2}
\]
For even $q$, division by two is modulo $A$. For odd $q$ the numerators
are even integers, and comparison modulo scalar signs is modulo $A/2$.
Changing a Spin lift changes only a scalar sign.
\end{lemma}
\begin{proof}
The Clifford construction on the exterior algebra of an isotropic
four-space gives the weights $\prod z_i^{s_i/2}$, $s_i=\pm1$, when the
natural eigenvalues are $z_i,z_i^{-1}$. The two parities of the number
of minus signs give the two half-spin modules. Indeed the $i$th torus
factor acts with weight $z_i^{1/2}$ on exterior monomials containing
the $i$th basis vector and with weight $z_i^{-1/2}$ on the others.
Conjugation on the creation and contraction operators then has the
required weights $z_i,z_i^{-1}$. Thus here the weights are the
half-sums of $1,q,q^2,ub$ in the two parities.

Use $q^3=-1$ and $qb=-b$ modulo $A$. Successively multiplying the
all-positive half-sum by $q$ gives six even-parity half-sums;
the other two are $\pm(1-u)b/2$. Starting with
$1+q+q^2-ub$ gives the odd-parity list. For odd $q$ the doubled
identities modulo $A$ prove the desired identities modulo $A/2$.
For even $q$, two is invertible modulo $A$.
\end{proof}

For $q\ge3$, the order of $q$ modulo $b$ is six.
It divides six because $q^3=-1$; $b>q+1>2$ excludes orders one
and three, and $q^2=1$ would force $q=2\pmod b$, excluding order two.
Moreover
\[
                              c_\epsilon\equiv q\pmod b.     \tag{3}
\]
Here $1+q+q^2\equiv2q\pmod b$ and $b$ is odd.

\begin{proposition}\label{generic}
Suppose $q$ is even and $q\ge4$, or $q$ is odd and $q\ge7$.
Choose a unit $u\pmod B$ other than $1$ and $-1$.
Every involutory automorphism of $G$ inverting $x$ belongs to $J$.
There are at most $AB$ such automorphisms.
Furthermore, $x$ belongs to no triality image of a stabilizer of a
nonsingular line in a natural eight-dimensional module.
\end{proposition}
\begin{proof}
For even $q$ take $u=2$. For odd $q\ge7$, $B$ is even and at least
eight, and $\varphi(B)>2$. The prime-power formula for Euler's function
shows that $\varphi(n)=2$ only for $n=3,4,6$.

Put $M=A$ for even $q$ and $M=A/2$ for odd $q$. By (3), the six
exponents $c_\epsilon q^j$ have period six already modulo $b$;
the two exponents $\pm d_\epsilon$ are divisible by $b$ and cannot
join that orbit. A similarity up to scalar signs between a half-spin
and the natural spectrum must therefore match their six-element orbits.
It implies $c_\epsilon=q^j\pmod M$ for some $j$. Reduction modulo $b$
forces $j=1\pmod6$. Since
\[
                       2(c_\epsilon-q)=(1+\epsilon u)b,
\]
the congruence implies $u=-\epsilon\pmod B$, a contradiction.
Taking inverses does not change either spectrum, so this also excludes
a diagram transposition moving the natural module for an inverter.

A nontrivial involutory field projection has $q=p^f$ with even $f$
and exponent $f/2$. If its diagram part moves the natural module,
the same comparison, reduced modulo $b$, gives
$p^{f/2}q=q^j$ after absorbing inversion into the six-element orbit.
Squaring gives $q^3=q^{2j}$, impossible for an element of order six.
Ordinary natural field inverters were excluded in the acquired torus
proof. An involution cannot project to a three-cycle in $S_3$.

The remaining diagram transposition fixes the natural module.
Apply $q^3$-Frobenius on the six-space and $q$-Frobenius on the
two-space. This is an involutory isometry inverting $X$.
It is proper: its determinant in odd characteristic is
$(-1)^3(-1)=1$, and its Dickson invariant in characteristic two is
$3+1=0\pmod2$. It therefore represents an element of $J$.
Any other natural inverter differs from it by an element of the
natural projective centralizer, the connected torus $T$.
Thus no improper natural inverter exists. All inverters lie in $J$,
where their pairwise quotients centralize $x$, proving the bound $|T|$.

Finally the natural module has no scalar-sign eigenvalue. In either
half-spin module a large-orbit exponent is nonzero modulo $b$.
A small exponent is zero modulo $M$ only if $u=\epsilon\pmod B$,
also excluded. There is no scalar-sign eigenvalue in any module.
A nonsingular line fixed by a projective isometry has eigenvalue
$1$ or $-1$ on an isometric lift, since its nonzero quadratic value
is preserved. Such a line therefore does not exist.
\end{proof}

\begin{proposition}\label{five}
For $q=5$ take $u=1$. At most $2AB$ involutory automorphisms invert
$x$. Their fixed subgroups in $G$ each have at most $24q^8$ elements
of square one, including the identity. There are at most six triality
images of nonsingular-line stabilizers containing $x$.
\end{proposition}
\begin{proof}
Here $A=126$, $B=6$, $b=21$, and $M=63$.
Formula (2) gives one half-spin spectrum identical modulo signs to
the natural one. The other has large-orbit representative $26$, of
order $63$, and a scalar two-space. The natural and first half-spin
modules have no scalar-sign eigenvalue, so the remaining module is
intrinsically distinguished. An inverter must fix its diagram point.
There are at most two possible diagram projections and no field
automorphism. Each such coset of $J$ has at most $|T|$ inverters.

Use the distinguished module to represent an inverter by a conformal
map $a$. Its six-space has irreducible eigenvalue order $63$.
Projective inversion has scalar $1$: the degree argument forces
the field exponent to be three, already giving exact inversion.
On that block, the square and conformal multiplier of $a$ are the
same scalar $c$, as in the natural torus proof. Since its projective
action is involutory, $a^2=cI$ on the entire module.

If $c$ is a square, the six-space contributes eigenspace dimensions
three and three. The two-space adds $2,0$, $1,1$, or $0,2$.
Thus the full pair is $(5,3)$, $(4,4)$, or $(3,5)$.
The written odd-characteristic projective fixed-group bounds give
$18q^8$ in the unequal cases and $24q^8$ in the balanced case.
If $c$ is nonsquare, the two eigenspaces over the algebraic closure
each have dimension four. The same bound applies, including the
quadratic-field form and projective kernel.

Only the distinguished module can have a fixed nonsingular line.
Its scalar two-space has minus type: the irreducible six-space is
negative and the eight-space is positive. Its $q+1=6$ lines are all
nonsingular, and they are the only possible lines. There is at most
one stabilizer for each line. This counts actual stabilizers, without
replacing their multiplicity by the number of conjugacy types.
\end{proof}

\paragraph{Attribution and limits.}
The three representations and triality action, finite automorphism
structure and Clifford construction are prior theory. The exponent
calculation, primitive choice separating spectra, and resulting
coset and actual-line bounds are the arguments here.
The cases $q=2,3$ have acquired finite witnesses.
The separate maximal-overgroup cover and final generation inequality
remain necessary; no whole-original or new-conjecture claim is made.

\begin{thebibliography}{9}\small
\bibitem{BT} D.~Barry and J.-P.~Tignol, \emph{Trialitarian triples},
arXiv:2301.10130v1, Sections~3.1--3.4.
This auxiliary source does not enlarge the fixed original catalogue.
\bibitem{BSW} D.~Bubboloni, P.~Spiga and Th.~Weigel,
\emph{Normal $2$-coverings of the finite simple groups and their
generalizations}, arXiv:2208.08756v1, Section~8, physical pages~81--82
for the triality-related orthogonal line stabilizers.
\end{thebibliography}
\end{document}

